\documentclass[11pt]{article}

\usepackage[utf8]{inputenc}
\usepackage[T1]{fontenc}
\usepackage{lmodern}
\usepackage[margin=1in]{geometry}
\usepackage{amsmath,amssymb,amsthm}
\usepackage{booktabs}
\usepackage{microtype}
\usepackage[hidelinks]{hyperref}
\usepackage{cleveref}

% Run-in headings italic, never bold.
\makeatletter
\renewcommand\paragraph{\@startsection{paragraph}{4}{\z@}%
  {3.25ex \@plus1ex \@minus.2ex}{-1em}%
  {\normalfont\normalsize\itshape}}
\makeatother

\title{A Schur Bridge from Nested Clustered Optimization to Global Minimum
Variance}
\author{Peter Cotton\thanks{Microprediction.
Email: \texttt{peter.cotton@microprediction.com}. A companion
verification script is \texttt{paper/verify\_schur\_nco\_bridge.py}
at \texttt{github.com/microprediction/schur}.}}
\date{}

\newcommand{\R}{\mathbb{R}}
\newcommand{\ones}{\mathbf{1}}
\DeclareMathOperator{\diag}{diag}
\DeclareMathOperator{\Cov}{Cov}
\DeclareMathOperator{\Var}{Var}

\newtheorem{proposition}{Proposition}
\newtheorem{remark}{Remark}
\newtheorem{definition}{Definition}
\newtheorem{corollary}{Corollary}

\begin{document}
\maketitle

\begin{abstract}
Nested clustered optimization allocates within each cluster from the
cluster's own covariance block and then across the resulting cluster
portfolios. Block inversion says the unconstrained minimum-variance
portfolio has the same two-tier shape, with each block replaced by its Schur
complement against every other asset. Conditioning instead on one knot from each other cluster truncates the
conditioning set in the manner of a Vecchia approximation, and we give the rank-one model of cross-cluster dependence under which it is
exact. Damping the complement by $\gamma\in[0,1]$ then gives a bridge with
nested clustered optimization at $\gamma=0$ and the global optimum at
$\gamma=1$, with no linear solve larger than a cluster or the number of
clusters. Under estimation error the optimal $\gamma$ can be strictly
interior and full coupling can remain locally optimal, and we give the local
theorem at the minimum-variance end with exact examples of both.
\end{abstract}

\section{Introduction}

Two-stage cluster methods partition the assets, allocate inside each
cluster, and then allocate across clusters. In nested clustered
optimization, NCO \cite{lopezdeprado2020}, the inner step optimizes
each cluster on its own covariance block and the outer step optimizes across
the resulting cluster portfolios. The outer problem has one dimension per
cluster, which is the source of the method's numerical stability.

The inner step reads only the cluster's own block. The only channel by which
one cluster's weights respond to another cluster is its scalar budget.

Schur-complementary allocation \cite{cotton2024schur} removes the same
blind spot from hierarchical risk parity \cite{lopezdeprado2016}. At
each split of a bisection tree it replaces a block by its Schur complement
against the sibling block, damped by $\gamma\in[0,1]$, and carries a
companion vector alongside. The result is a bridge with hierarchical risk
parity at $\gamma=0$ and minimum variance at $\gamma=1$.

This note builds the corresponding bridge for a flat partition. The left
end is NCO and the right end is the unconstrained global minimum-variance
portfolio. The obstacle is cost, since the complement of a cluster against
everything else needs an inverse of nearly full size.

Spatial statistics meets the same obstacle by conditioning on a few chosen
indices rather than on everything \cite{vecchia1988}. The conditionals
there are Schur complements, so the device transfers directly.

We borrow its vocabulary. One member of each cluster is its \emph{knot}, as
in the low-rank spatial models of Banerjee et al.~\cite{banerjee2008} that
Katzfuss and Guinness~\cite{katzfuss2021} place inside the Vecchia framework. The other
clusters' knots are a cluster's \emph{conditioning set}. We state the model
under which conditioning on that set loses nothing.

\Cref{sec:where} asks where to sit on the bridge when the covariance is
estimated. Part of the analysis of the tree in \cite{cotton2026bridge}
transfers, and exact examples show that both an interior optimum and a
locally optimal full coupling occur.

All exactness statements concern the unconstrained, fully invested,
long-short problem. With box or long-only constraints no block
decomposition of the optimum is exact, here or on the tree.

\section{Setup}

Let $\Sigma\succ0$ be the covariance of $n$ assets partitioned into
clusters $I_1,\dots,I_k$. Cluster $i$ has a knot $p_i\in I_i$ and
remaining members $J_i=I_i\setminus\{p_i\}$. Write $P=\{p_1,\dots,p_k\}$
and $P_{-i}=P\setminus\{p_i\}$. For index sets $X,Y$ write $\Sigma_{XY}$ for the
corresponding block, $\Sigma_{ii}$ for the block of cluster $i$, and $-i$
for the complement of $I_i$.

A two-stage method chooses a portfolio $v_i$ for each cluster, supported on
$I_i$ with $\ones^\top v_i=1$, and then budgets across the clusters. With
$V$ the $n\times k$ matrix of cluster portfolios,
\begin{equation}
  a \;=\; \operatorname{alloc}\!\big(V^\top\Sigma V\big),
  \qquad w \;=\; Va .
  \label{eq:twostage}
\end{equation}
NCO takes $v_i$ to be the minimum-variance portfolio of $\Sigma_{ii}$ and
minimum variance for the outer allocation.

\section{Block inversion}

Fix a vector $u$ and a cluster $i$. Block inversion of $\Sigma$ on the
partition $(I_i,-i)$ gives \cite[App.~A]{cotton2024schur}
\begin{equation}
  \big(\Sigma^{-1}u\big)_{I_i}
  \;=\; \big(\Sigma^{c}_{ii}\big)^{-1} b_i,
  \label{eq:blockinv}
\end{equation}
where
\begin{equation}
  \Sigma^{c}_{ii} = \Sigma_{ii}-\Sigma_{i,-i}\Sigma_{-i,-i}^{-1}\Sigma_{-i,i},
  \qquad
  b_i = u_{I_i}-\Sigma_{i,-i}\Sigma_{-i,-i}^{-1}u_{-i} .
  \label{eq:complement}
\end{equation}
For a positive-definite $Q$ and a vector $b$ we call $Q^{-1}b$ a
\emph{generalized minimum-variance direction}. Up to scale it minimizes
$x^\top Qx$ subject to $b^\top x=1$, which is ordinary minimum variance only
when $b=\ones$.

With $u=\ones$, the blocks of the global minimum-variance direction
$\Sigma^{-1}\ones$ are therefore generalized minimum-variance directions,
one per cluster, on the cluster's covariance conditional on everything
else. There is no cross-cluster optimizer in the global optimum. The
cluster totals are whatever the stacked directions sum to, and every piece
of cross-cluster information enters through the conditioning.

NCO has the same shape with the conditioning removed. Restoring it costs an
inverse of size $n-|I_i|$ per cluster, which is what a two-stage method
exists to avoid.

\section{The gateway model}

\begin{definition}[Gateway model]
For every cluster $j$, the residual of the remaining members $J_j$ after linear
regression on their own knot is uncorrelated with every asset outside
$I_j$. Equivalently the cross-cluster residual
\begin{equation}
  R_j(p_j) \;=\; \Sigma_{J_j,-j}-\Sigma_{J_j p_j}\,\sigma_{p_j}^{-2}\,\Sigma_{p_j,-j}
  \label{eq:residual}
\end{equation}
is zero.
\end{definition}

A member may be correlated with anything, but its correlation with other
clusters passes entirely through its knot. The knot is the cluster's only
gateway. Each cross-cluster block of $\Sigma$ then has rank at most one.

This is a strong assumption, and clustering alone does not imply it. It
asks that member residuals be uncorrelated both with the other knots and
with the other clusters' member residuals.

\begin{proposition}[Knots are sufficient]
\label{prop:sufficient}
Under the gateway model, for each cluster $i$ and every $u$,
\begin{equation}
  \Sigma^{c}_{ii}
  \;=\; \Sigma_{ii}-\Sigma_{i,P_{-i}}\,\Sigma_{P_{-i}P_{-i}}^{-1}\,\Sigma_{P_{-i},i},
  \qquad
  b_i \;=\; u_{I_i}-\Sigma_{i,P_{-i}}\,\Sigma_{P_{-i}P_{-i}}^{-1}\,u_{P_{-i}} .
  \label{eq:cheap}
\end{equation}
\end{proposition}

\begin{proof}
The matrix $\Sigma_{i,-i}\Sigma_{-i,-i}^{-1}$ holds the coefficients of the
linear regression of cluster $i$ on all other assets. Split the regressors
into the other knots $P_{-i}$ and the other remaining members $J_{-i}$, and write
$J_{-i}=\beta\,P_{-i}+\varepsilon$ for the block-diagonal regression of
each set of remaining members on its own knot. The map from $(P_{-i},J_{-i})$ to
$(P_{-i},\varepsilon)$ is invertible, so the regression may be run on the
latter pair. Under the gateway model $\varepsilon$ is uncorrelated with
cluster $i$ and with $P_{-i}$, so its coefficients vanish and the
coefficients on $P_{-i}$ are $\Sigma_{i,P_{-i}}\Sigma_{P_{-i}P_{-i}}^{-1}$.
Substituting into \eqref{eq:complement} gives \eqref{eq:cheap}, with
$u_{-i}$ contributing only through its $P_{-i}$ entries.
\end{proof}

Conditioning cluster $i$ on the other knots is the same as conditioning it
on everything. The inverse in \eqref{eq:cheap} is $(k-1)\times(k-1)$.

\begin{remark}[Relation to Vecchia approximations]
A Vecchia approximation orders the variables and replaces each conditioning
set by a small subset of the predecessors, which defines one acyclic
approximate joint law \cite{vecchia1988,katzfuss2021}. The conditionals
are Schur complements. Replacing $-i$ by $P_{-i}$ in \eqref{eq:complement}
is the same truncation of a conditioning set, but it is applied to the full
conditional of every cluster at once. Those conditionals are mutually
cyclic and no approximate joint law is formed, so the recipe is
Vecchia-inspired rather than a Vecchia approximation.
\end{remark}

The gateway model itself does have an exact Vecchia form.

\begin{corollary}[Gaussian factorization]
\label{cor:gaussian}
If the returns $r$ are jointly Gaussian, the gateway model holds if and only
if
\begin{equation}
  f(r) \;=\; f(r_P)\,\prod_{i=1}^{k} f\big(r_{J_i}\mid r_{p_i}\big).
  \label{eq:dag}
\end{equation}
\end{corollary}

\begin{proof}
Under the model the residuals $\varepsilon_i=r_{J_i}-\beta_ir_{p_i}$ are
uncorrelated with $r_P$ and with each other, hence independent under
Gaussianity, which gives \eqref{eq:dag}. Conversely \eqref{eq:dag} makes
$\varepsilon_i$ independent of every asset outside $I_i$.
\end{proof}

With the knots ordered first, \eqref{eq:dag} is a block-Vecchia
factorization, a sparser relative of the knot-based cases collected by
Katzfuss and Guinness~\cite{katzfuss2021}. The covariance results in this note do not need
Gaussianity.

\begin{remark}[Testing the model]
Regressing each remaining member on its own knot and the other knots, and
checking that the other-knot coefficients vanish, does not test the model.
Member residuals can be orthogonal to every knot and still be correlated
with each other across clusters, and the sufficiency of the knots then
fails. The model holds exactly when $R_j(p_j)=0$ for every $j$, which can
be checked from the sample covariance before any portfolio is built.
\end{remark}

\begin{remark}[Measuring the loss and choosing knots]
The raw size of $R_j$ does not measure what is lost when the model fails. A
residual with variance $\delta^2$ and covariance $c\delta$ with an outside
asset has $R=c\delta\to0$, while the Schur correction it carries, $c^2$,
does not move. The loss for cluster $i$ is the omitted correction
\begin{equation}
  \Delta_i \;=\; \Sigma_{i,-i}\Sigma_{-i,-i}^{-1}\Sigma_{-i,i}
   -\Sigma_{i,P_{-i}}\Sigma_{P_{-i}P_{-i}}^{-1}\Sigma_{P_{-i},i}
  \;\succeq\; 0,
  \label{eq:loss}
\end{equation}
the part of cluster $i$ explained by the other assets beyond their knots.
A scale-free screen for a candidate knot $p\in I_j$ is the largest canonical
correlation between the residual of $I_j\setminus\{p\}$ on $p$ and the
assets outside $I_j$. Both need a large inverse, but they are diagnostics
computed once and are not part of the allocation.
\end{remark}

\begin{remark}[A symmetric factor model]
Suppose instead that every member of cluster $j$ loads on a latent cluster
factor, $r_m=\lambda_m f_j+\eta_m$, with idiosyncratic terms uncorrelated
with each other and with the factors. Then
$R_j(p)=(1-\rho_p)\,\lambda_{J_j}\Cov(f_j,r_{-j})$, where
$\rho_p=\lambda_p^2\Var(f_j)/\sigma_p^2$ is the share of the knot's variance
explained by the factor. The gateway model holds to the extent the knot is
the factor. A knot with a high $\rho_p$ is then the natural choice, and the
member most aligned with the cluster's first principal component
approximates it.
\end{remark}

\section{The bridge}

Following \cite{cotton2024schur}, damp the conditioning. For
$\gamma\in[0,1]$ define the \emph{conditioned pair} of cluster $i$,
\begin{equation}
\begin{aligned}
  Q_i(\gamma) &= \Sigma_{ii}-\gamma\,\Sigma_{i,P_{-i}}\Sigma_{P_{-i}P_{-i}}^{-1}\Sigma_{P_{-i},i},\\
  b_i(\gamma) &= u_{I_i}-\gamma\,\Sigma_{i,P_{-i}}\Sigma_{P_{-i}P_{-i}}^{-1}u_{P_{-i}},
\end{aligned}
  \label{eq:pair}
\end{equation}
and the direction $d_i(\gamma)=Q_i(\gamma)^{-1}b_i(\gamma)$. Let $D_\gamma$
be the $n\times k$ matrix with $d_i(\gamma)$ in column $i$, supported on
$I_i$. The bridge portfolio is
\begin{equation}
  w_\gamma \;\propto\; D_\gamma\big(D_\gamma^\top\Sigma D_\gamma\big)^{-1}D_\gamma^\top u .
  \label{eq:top}
\end{equation}
It solves, on the span of the $k$ directions, the problem that
$\Sigma^{-1}u$ solves on the whole space. A direction that vanishes is
dropped.

\paragraph{Two-tier reading.} The right side of \eqref{eq:top} does not
change when a column of $D_\gamma$ is rescaled. When every total
$\ones^\top d_i(\gamma)$ is nonzero, set $v_i=d_i/\ones^\top d_i$. Then
\eqref{eq:top} is the two-stage method \eqref{eq:twostage} with
$a\propto(V^\top\Sigma V)^{-1}V^\top u$. The induced outer inputs are the
$k\times k$ covariance $V^\top\Sigma V$ and the $k$-vector $V^\top u$, which
is $\ones$ when $u=\ones$. The unnormalized form \eqref{eq:top} has no
singularity when a cluster total crosses zero.

\begin{proposition}[The two ends of the bridge]
\label{prop:bridge}
\leavevmode
\begin{enumerate}
\item At $\gamma=0$ with $u=\ones$ the recipe is NCO with minimum variance
  at both tiers.
\item Under the gateway model the columns of $D_1$ are the blocks of
  $\Sigma^{-1}u$, and $w_1\propto\Sigma^{-1}u$.
\end{enumerate}
\end{proposition}

\begin{proof}
At $\gamma=0$ the pair is $(\Sigma_{ii},\ones)$, so $d_i=\Sigma_{ii}^{-1}\ones$
has a positive total, $v_i$ is the minimum-variance portfolio of the block,
and the outer tier is minimum variance across the cluster portfolios. For
the second claim, \eqref{eq:blockinv} and \cref{prop:sufficient} give
$\Sigma^{-1}u=D_1\ones$, which lies in the column span of $D_1$. Then
$D_1(D_1^\top\Sigma D_1)^{-1}D_1^\top u
=D_1(D_1^\top\Sigma D_1)^{-1}D_1^\top\Sigma D_1\ones=D_1\ones$.
\end{proof}

\begin{table}[ht]
\centering
\begin{tabular}{@{}ll@{}}
\toprule
$\gamma$ & Method \\
\midrule
$0$ & NCO \cite{lopezdeprado2020} \\
$(0,1)$ & the interior of the bridge \\
$1$ & unconstrained global minimum variance, under the gateway model \\
\bottomrule
\end{tabular}
\caption{The bridge with $u=\ones$ and minimum variance at both tiers.}
\label{tab:bridge}
\end{table}

\paragraph{Scope.} The right end is the unconstrained, fully invested,
long-short solution. No condition on cluster totals is needed for
\eqref{eq:top}. A global solution can hold a self-financing position inside
a cluster with zero net weight. The two-tier reading fails there, because
no cluster budget times a normalized cluster portfolio represents it, and
\eqref{eq:top} does not.

\begin{remark}[Other objectives]
\Cref{prop:bridge} is stated for a direction $\Sigma^{-1}u$. With $u=\mu$
the direction $\Sigma^{-1}\mu$ maximizes the Sharpe ratio. Its fully
invested multiple does so only when $\ones^\top\Sigma^{-1}\mu>0$. When that
total is negative the fully invested multiple has the opposite orientation
and the lowest signed Sharpe ratio.

Mean--variance problems with a budget
constraint have solutions that combine $\Sigma^{-1}\mu$ and
$\Sigma^{-1}\ones$, and each part is covered separately. Risk parity is not
of this form. The recipe still runs with a risk-parity outer tier, but
nothing is exact at $\gamma=1$.
\end{remark}

\begin{remark}[Rules that read only a covariance]
An inner rule that cannot take a constraint vector can be applied after a
change of variables. With $y=b_i\circ x$ the constraint $b_i^\top x=1$
becomes $\ones^\top y=1$ and the covariance becomes
$M=\diag(b_i)^{-1}Q_i\diag(b_i)^{-1}$. Apply the rule to $M$ and map back by
$x=y/b_i$, coordinate by coordinate. For minimum variance this returns
$Q_i^{-1}b_i$. It requires $b_i$ to have no zero entry. Applying minimum
variance to $M$ without mapping back returns $\diag(b_i)Q_i^{-1}b_i$, which
is a different portfolio.
\end{remark}

\begin{remark}[Relation to the tree]
Schur complements compose: the complement of a complement is the complement
against the union. So at $\gamma=1$ the recursion of
\cite{cotton2024schur}, run down any tree whose leaves are the clusters
and stopped there, hands each cluster the pair \eqref{eq:complement}. The
flat construction and the iterated one agree at the right end. They differ
in the interior, because the tree damps at every split and the damping
compounds along the path, while \eqref{eq:pair} damps each cluster once
against all the other knots.
\end{remark}

\section{Where to sit on the bridge}
\label{sec:where}

The cross-knot regression coefficients in \eqref{eq:pair} are estimated, and
$\gamma$ is how much of that estimate to trust. The analysis of the tree in
\cite{cotton2026bridge} transfers in part. This section takes $u=\ones$.

Let $\widehat\Sigma_\tau=\Sigma+\tau E$ be an estimate, where $E$ is a random
symmetric matrix whose law is unchanged by $E\mapsto-E$, and
$\widehat\Sigma_\tau\succ0$ on the range of $\tau$ considered. Every
quantity in the recipe is computed from $\widehat\Sigma_\tau$. The partition
and the knots are held fixed as $\tau$ varies, since reclustering or
reselecting knots from the data can destroy smoothness. The portfolio
$w_\gamma$ of \eqref{eq:top} is normalized by $\ones^\top w_\gamma=1$, and
we assume the normalizing total stays nonzero. Define the expected
out-of-sample variance and its noiseless profile,
\begin{equation}
  F(\gamma,\tau)=\mathbb{E}\big[w_\gamma(\widehat\Sigma_\tau)^\top\,\Sigma\,
  w_\gamma(\widehat\Sigma_\tau)\big],
  \qquad V_0(\gamma)=F(\gamma,0).
  \label{eq:F}
\end{equation}
It is even in $\tau$. We assume $F$ is smooth, put
$G(\gamma)=\tfrac12\,\partial_\tau^2F(\gamma,0)$, and assume the expansion
$F=V_0+\tau^2G+O(\tau^4)$ holds in $C^2$ as a function of $\gamma$.

\begin{proposition}[Incremental noise cost]
\label{prop:incremental}
Let $\Phi(\gamma,\tau)=F(\gamma,\tau)-F(0,\tau)-\big[V_0(\gamma)-V_0(0)\big]$.
Then $\Phi(\gamma,\tau)=\gamma\,\tau^2H(\gamma,\tau)$ for a smooth $H$.
\end{proposition}

\begin{proof}
$\Phi(0,\tau)=0$, so $\Phi=\gamma\,\Phi_1$ with $\Phi_1$ smooth. Also
$\Phi(\gamma,0)=0$, so $\Phi_1(\gamma,0)=0$, and $\Phi_1$ is even in $\tau$,
so $\Phi_1=\tau^2H$.
\end{proof}

On the tree, under noise supported only on the top-level cross-block, the
hierarchical end does not read that block, and the uncentred cost $F-V_0$
already carries the factor $\gamma$ \cite{cotton2026bridge}. Under
whole-matrix noise that fails on the tree as well, because $F(0,\tau)$ is
then noisy.

NCO is different even under cross-cluster noise alone. Its outer tier reads
cross-cluster covariance at $\gamma=0$, so NCO has an estimation cost of its
own, and the factor $\gamma$ appears only after centring at the noisy NCO
end. The centred statement holds for the general noise of this section, and
the cost of moving away from NCO is still of order $\gamma\tau^2$.

\begin{corollary}[Separation from NCO]
\label{cor:separation}
If $V_0'(0)<0$ then $\partial_\gamma F(0,\tau)<0$ for all sufficiently small
$\tau$, and NCO is strictly suboptimal.
\end{corollary}

\begin{proof}
By \cref{prop:incremental},
$\partial_\gamma F(0,\tau)=V_0'(0)+\tau^2H(0,\tau)$.
\end{proof}

\begin{proposition}[The minimum-variance end]
\label{prop:mvend}
Under the gateway model $V_0'(1)=0$ and
$V_0''(1)=2\,\dot w^\top\Sigma\,\dot w$, where
$\dot w=\partial_\gamma w_\gamma$ at $\gamma=1$. Suppose $\dot w\neq0$. Then
for all sufficiently small $\tau$ the map $F(\cdot,\tau)$ has a unique
stationary point near $1$, a strict local minimizer, at
\begin{equation}
  \tilde\gamma(\tau)\;=\;1-\frac{G'(1)}{V_0''(1)}\,\tau^2+O(\tau^4).
  \label{eq:shift}
\end{equation}
If $G'(1)>0$ this local minimizer is strictly inside the bridge for small
$\tau>0$. If $G'(1)<0$ then $\gamma=1$ is a strict local minimizer of
$F(\cdot,\tau)$ on $[0,1]$. If $G'(1)=0$ the second order is silent and the
fourth or a higher order decides.

Suppose in addition that $\gamma=1$ is the unique global minimizer of $V_0$
on $[0,1]$ and that $F(\cdot,\tau)\to V_0$ uniformly on $[0,1]$. Then for
all sufficiently small $\tau$ these local minimizers are the global optimum
over $[0,1]$.
\end{proposition}

\begin{proof}
Since $\ones^\top w_\gamma=1$ we have $\ones^\top\dot w=0$ and
$\ones^\top\ddot w=0$. By \cref{prop:bridge},
$w_1=\Sigma^{-1}\ones/\ones^\top\Sigma^{-1}\ones$, so
$\Sigma w_1\propto\ones$. Hence $V_0'(1)=2\dot w^\top\Sigma w_1=0$, and in
$V_0''(1)=2\dot w^\top\Sigma\dot w+2\ddot w^\top\Sigma w_1$ the second term
vanishes. Apply the implicit function theorem to
$\partial_\gamma F(\gamma,\tau)=0$ at $(1,0)$, where
$\partial_\gamma^2F=V_0''(1)>0$. This gives \eqref{eq:shift} and strict
convexity of $F(\cdot,\tau)$ on a neighbourhood of $1$ that does not depend
on $\tau$.

When $G'(1)<0$ the stationary point lies beyond $1$, so
$F(\cdot,\tau)$ decreases up to the boundary. Under the additional
hypotheses the global minimizers of $F(\cdot,\tau)$ converge to $1$, so for
small $\tau$ they lie in that neighbourhood.
\end{proof}

Gateway exactness makes $w_1$ the minimum-variance portfolio. It does not by
itself exclude another $\gamma<1$ that produces the same portfolio, which is
why uniqueness is a hypothesis of the global statement.

The structure of the bridge does not determine the sign of $G'(1)$. Both
signs occur, and each is strict in the examples below, so each persists on
a relatively open set within the gateway-model parameterization.

\paragraph{An exact interior optimum.} Take clusters $\{1,2\}$ and
$\{3,4\}$ with knots $1$ and $3$, and
\begin{equation*}
  \Sigma=\begin{pmatrix}1&0&\tfrac12&0\\0&1&0&0\\\tfrac12&0&4&0\\0&0&0&1\end{pmatrix}.
\end{equation*}
The estimate replaces the cross-knot entry $\tfrac12$ by $Z\in\{0,1\}$, each
with probability $\tfrac12$. The population and both estimates are positive
definite and satisfy the gateway model, because the remaining members are
uncorrelated with everything.

On the branch $Z=0$ the directions do not depend on $\gamma$ and the
population variance is $56/169$. On the branch $Z=1$ the directions are
$d_1=(1,1)$ and $d_2=(x,1)$ with $x=(1-\gamma)/(4-\gamma)\in[0,\tfrac14]$,
and the population variance of the resulting portfolio is
\begin{equation*}
  V(x)=\frac{49x^4-7x^3+36x^2-2x+6}{2\,(7x^2+3)^2},
  \qquad
  V'(x)=\frac{P(x)}{2\,(7x^2+3)^3},
\end{equation*}
with $P(x)=49x^4+84x^3-21x^2+48x-6$. On $[0,\tfrac14]$ we have
$P'(x)=196x^3+252x^2-42x+48>0$, while $P(0)=-6$ and $P(\tfrac14)=1585/256$.
So $P$ has exactly one root $x_\star\approx0.12823$ there, and
\begin{equation*}
  \gamma_\star=\frac{1-4x_\star}{1-x_\star}\approx0.5587
\end{equation*}
is the unique optimum, strictly inside the bridge. In exact arithmetic
$F(0)-F(\tfrac12)=3/1210$ and $F(1)-F(\tfrac12)=5/1452$.

The same covariance with small symmetric noise, $Z=\tfrac12\pm\tau$,
illustrates \cref{prop:incremental,prop:mvend}. There
$G'(1)=31779128936/554523204167\approx0.0573$ and
$V_0''(1)=48136/3571279\approx0.01348$ exactly, so \eqref{eq:shift}
predicts $\tilde\gamma(0.05)\approx0.9894$ against a computed optimum of
$0.9896$. The function $H$ of \cref{prop:incremental} is close to $0.0071$
at every $\gamma,\tau\in\{0.05,0.2\}$.

\paragraph{Full coupling stays locally optimal.} Take three clusters, each a
knot and one remaining member that is uncorrelated with everything, with
member variances $1$, $4$ and $8$. Let the knot covariance and its estimates
be
\begin{equation*}
  S=\begin{pmatrix}5&6&5\\6&\tfrac{33}{2}&12\\5&12&\tfrac{25}{2}\end{pmatrix},
  \qquad
  \widehat S_\tau=S\pm\tau\begin{pmatrix}0&-1&0\\-1&0&2\\0&2&0\end{pmatrix},
\end{equation*}
the two signs equally likely. The noise is symmetric, touches only
cross-knot entries, and preserves the gateway model. Exact differentiation
gives $V_0''(1)\approx2.3949$ and
\begin{equation*}
  G'(1)=-\frac{4720238484060245648386156800}{806907939268294475102923640251}
  \approx-0.0058498<0,
\end{equation*}
so by \cref{prop:mvend} full coupling is a strict local minimizer for all
sufficiently small $\tau$.

As a spot check, at $\tau=\tfrac1{10}$ and $\tau=\tfrac1{100}$ exact
evaluation gives $F(\gamma,\tau)>F(1,\tau)$ for
$\gamma\in\{0.999,\,0.99,\,0.9,\,0.5,\,0\}$. This is consistent with a
global optimum at full coupling but does not prove one.

Neither end owns the bridge. An interior optimum is not a consequence of the
architecture, and a statement that it is generic needs a noise model. The
derivatives quoted in both examples are exact, computed by truncated
bivariate jets in rational arithmetic.

\section{The outer covariance}

The gateway model also gives the outer covariance a simple form.

\begin{proposition}[Knot covariance plus a ridge]
\label{prop:outer}
Under the gateway model, let $v_i$ be any portfolios supported on their
clusters. Write $\beta_i=\Sigma_{J_ip_i}/\sigma_{p_i}^2$ for the member
betas, $E_i=\Sigma_{J_iJ_i}-\beta_i\beta_i^\top\sigma_{p_i}^2$ for the
residual covariance, and $\kappa_i=v_{i,p_i}+\beta_i^\top v_{i,J_i}$ for
the cluster portfolio's beta to its knot. Then
\begin{equation}
  V^\top\Sigma V \;=\; K\,\Sigma_{PP}\,K + D,
  \qquad K=\diag(\kappa_i),
  \quad D=\diag\!\big(v_{i,J_i}^\top E_i\,v_{i,J_i}\big).
  \label{eq:outer}
\end{equation}
\end{proposition}

\begin{proof}
Under the gateway model $\Sigma_{J_i,l}=\beta_i\Sigma_{p_i,l}$ for every
asset $l$ outside $I_i$. Applying this on both sides gives
$v_i^\top\Sigma_{I_iI_l}v_l=\kappa_i\kappa_l\Sigma_{p_ip_l}$ for $i\neq l$.
On the diagonal, the knot is uncorrelated with the member residual, so
$v_i^\top\Sigma_{ii}v_i=\kappa_i^2\sigma_{p_i}^2+v_{i,J_i}^\top E_iv_{i,J_i}$.
\end{proof}

The outer problem sees the knot covariance, scaled by the cluster
portfolios' knot betas, plus their residual variances on the diagonal. The
diagonal term raises the smallest eigenvalue of the outer covariance by at
least the smallest residual variance. The remaining members supply a ridge, which
is one reading of why the outer tier of NCO is stable.

\section{Cost}

No step forms an $n\times n$ inverse. The largest linear solve has dimension
$\max(k,\max_i|I_i|)$, as in NCO: the outer tier is $k\times k$ and each
conditioning inverse is $(k-1)\times(k-1)$. The bridge adds work that NCO
does not do: the cross-block products $\Sigma_{i,P_{-i}}$, and a
leave-one-knot-out inverse per cluster, obtainable from $\Sigma_{PP}^{-1}$
by a rank-one downdate.

If the gateway model fails, \eqref{eq:cheap} approximates the full
complement in \eqref{eq:complement}, and $\Delta_i$ in \eqref{eq:loss} is
what is lost.

\section{Closing}

The contribution is deliberately small. The identity is block inversion, the
truncation of the conditioning set follows Vecchia approximations, and the
damping is from
\cite{cotton2024schur}. What is added is the model under which one knot
per cluster suffices, and the bridge from NCO to the global optimum that it
makes cheap.

\Cref{prop:sufficient} may also serve Schur-complementary allocation on the
tree, where the complement against a sibling block could be taken against
that block's knots. The whole-matrix analysis of \cite{cotton2026bridge}
rests on a scalar reduction specific to the tree and needs a replacement
here. We leave both, and an empirical comparison, to future work.

\begin{thebibliography}{99}

\bibitem{lopezdeprado2020}
M.~L\'opez de Prado,
\emph{Machine Learning for Asset Managers},
Cambridge University Press, 2020.

\bibitem{cotton2024schur}
P.~Cotton,
\emph{Schur complementary allocation: a unification of hierarchical risk
parity and minimum variance portfolios},
arXiv:2411.05807, 2024.

\bibitem{lopezdeprado2016}
M.~L\'opez de Prado,
\emph{Building diversified portfolios that outperform out of sample},
Journal of Portfolio Management \textbf{42}(4) (2016), 59--69.

\bibitem{vecchia1988}
A.~V.~Vecchia,
\emph{Estimation and model identification for continuous spatial processes},
Journal of the Royal Statistical Society: Series B \textbf{50}(2) (1988),
297--312.

\bibitem{banerjee2008}
S.~Banerjee, A.~E.~Gelfand, A.~O.~Finley, and H.~Sang,
\emph{Gaussian predictive process models for large spatial data sets},
Journal of the Royal Statistical Society: Series B \textbf{70}(4) (2008),
825--848.

\bibitem{katzfuss2021}
M.~Katzfuss and J.~Guinness,
\emph{A general framework for Vecchia approximations of Gaussian processes},
Statistical Science \textbf{36}(1) (2021), 124--141.

\bibitem{cotton2026bridge}
P.~Cotton,
\emph{When the out-of-sample-optimal Schur portfolio lies between HRP and
minimum variance}, working paper, 2026.
\url{https://schur.microprediction.org/neither-end-of-the-bridge.pdf}

\end{thebibliography}

\end{document}
